Project Goal:
Build a glucose prediction model that uses the fewest meaningful predictors to accurately predict glucose levels keeping prediction error low
Optimization: accuracy versus sparsity (fewest predictors)


Insulin Codes: 
33 = Regular insulin dose                               [short-acting insulin (works soon after injection, often around meals)]
34 = NPH insulin dose                                   [intermediate-acting insulin (lasts longer than regular, and is used for background insulin plan)]
35 = UltraLente insulin dose                            [Long-acting insulin (Different role than regular insulin, as its not for immediate spikes)]
48 = Unspecified blood glucose measurement              [Glucose was measured, but not tied to meal context]
57 = Unspecified blood glucose measurement              [^]
58 = Pre-breakfast blood glucose measurement
59 = Post-breakfast blood glucose measurement
60 = Pre-lunch blood glucose measurement
61 = Post-lunch blood glucose measurement
62 = Pre-supper blood glucose measurement
63 = Post-supper blood glucose measurement
64 = Pre-snack blood glucose measurement
65 = Hypoglycemic symptoms                              [symptoms of low blood sugar]
66 = Typical meal ingestion
67 = More-than-usual meal ingestion
68 = Less-than-usual meal ingestion
69 = Typical exercise activity
70 = More-than-usual exercise activity
71 = Less-than-usual exercise activity
72 = Unspecified special event

___________________________________________________________________
We predict a patients pre-meal glucose level using prior insulin, meal, exercise, and previous glucose levels, while selecting a spare set of predictors

# Variables And Definitions:
y_i: glucose value for observation i
x_i: feature vector for observation i

beta: model coefficients
z_j: binary variable indicating whether feature j is selected
c_j: cost or weight for selecting feature j

n: number of training observations
p: number of candidate features
tau: maximum acceptable prediction error
M: sufficiently large constant linking beta_j and z_j

You can say:
y_i comes only from pre-meal glucose rows
x_i may include previous glucose, insulin doses, meal indicators, exercise indicators, Day, MinuteOfDay, and lag features
___________________________________________________________________

# Ideal Optimization Model:
The exact sparse feature-selection problem is:

minimize over beta, z:
sum from j=1 to p of c_j z_j

subject to
(1/n) ||y - X beta||_2^2 <= tau
-M z_j <= beta_j <= M z_j,   for j = 1, ..., p
z_j in {0,1},   for j = 1, ..., p

Meaning
The objective minimizes the number or total cost of selected predictors.
The MSE constraint forces the model to stay accurate enough.
The linking constraint says that if z_j = 0, then beta_j must be zero, so feature j is not used.

___________________________________________________________________
# Convexity of the Ideal Model:
This ideal model is non-convex because:
z_j are binary variables
the feature-selection structure is combinatorial
___________________________________________________________________
# Convex Relaxation: Weighted LASSO Model
minimize over beta:
(1 / 2n) ||y - X beta||_2^2 + lambda * sum from j=1 to p of c_j |beta_j|

minimize over beta:
(1 / 2n) ||y - X beta||_2^2
+ lambda * sum from j=1 to p of c_j |beta_j|
+ (rho / 2) ||beta||_2^2

Meaning
The first term measures prediction error
The L1 penalty encourages sparsity
The optional L2 penalty improves stability when features are correlated
___________________________________________________________________

# Convexity Of The Relaxed Model
This relaxed model is convex because:
the squared error term is convex
the L1 norm is convex
the L2 squared norm is convex

If rho > 0, the elastic net version is strongly convex, which helps give a unique solution.
___________________________________________________________________
# Feasible Set

If you write LASSO in constrained form, the model becomes:
minimize over beta:
(1 / 2n) ||y - X beta||_2^2

subject to
sum from j=1 to p of c_j |beta_j| <= s

This feasible set is a weighted L1 ball, which is:
convex
closed
bounded
___________________________________________________________________
# Why A Solution Exists

For the convex relaxed model:
the objective is continuous
the feasible set in constrained form is closed and bounded
therefore a minimizer exists
For the elastic net version:

the objective is coercive and convex
so a minimizer exists
if rho > 0, the solution is typically unique
___________________________________________________________________
# Closed-Form Baseline
LASSO does not usually have a full closed-form solution, but ridge regression does, so ridge is a good mathematical baselin
Ridge model:
minimize over beta:
(1 / 2n) ||y - X beta||_2^2 + (rho / 2) ||beta||_2^2

Closed-form solution:
beta_hat = (X^T X + n rho I)^(-1) X^T y
___________________________________________________________________
# Numerical Solution For The Sparse Model

For LASSO or elastic net, use coordinate descent.

A standard update is:
beta_j <- S(a_j, lambda c_j) / b_j
where
a_j = (1/n) x_j^T (y - X_{-j} beta_{-j})
and
b_j = (1/n) ||x_j||_2^2 + rho
and the soft-thresholding operator is
S(a,t) = sign(a) * max(|a| - t, 0)
___________________________________________________________________
# Validation Plan

Use this as your workflow:
Define the target
    Use rows with glucose codes 58, 60, 62, 64.

Build features
    Use only events that happened before each target measurement.

Clean the data
    Handle malformed dates, irregular times, and nonnumeric glucose values if needed.

Split by time
    Train on earlier observations and test on later observations.

Compare models
    Mean predictor
    Ridge regression
    LASSO
    Elastic net
    Optional PCA baseline

Evaluate
    MSE
    RMSE
    Number of selected predictors
    tradeoff between sparsity and accuracy
___________________________________________________________________
Workflow

1. Load data
    Read all data-01 to data-70 patient files and combine them into one table.

2. Transform raw records
Convert:
    Date to patient-relative Day
    Time to MinuteOfDay
    Code to readable event labels

3. Choose target rows
Keep glucose rows with pre-meal codes:
    58 pre-breakfast
    60 pre-lunch
    62 pre-supper
    64 pre-snack

4. Build predictors for each target row
For each target glucose reading, create features using only earlier events:
    previous glucose
    previous insulin doses
    meal indicators
    exercise indicators
    Day
    MinuteOfDay
    time since last glucose reading
    optional lag summaries like recent average glucose

5. Clean and filter data
Handle:
    malformed dates
    nonnumeric glucose entries like 0Hi, 0Lo
    missing values
    rare or noisy codes if needed

6. Create train/test split
Split by time, not randomly:
    train on earlier records
    validate/test on later records
    This is better because we are predicting future glucose from past events.

# Models / Algorithms
7. Baseline 1: simple mean predictor
    Predict glucose using the average training glucose.
    This gives a weak baseline to beat.

8. Baseline 2: previous glucose predictor
    Predict the next pre-meal glucose using the most recent previous glucose.
    This is often a surprisingly strong baseline.

9. Ridge regression baseline
Fit ridge regression on all engineered features.
Purpose:
    closed-form mathematical baseline
    handles correlated predictors
    good comparison point

10. LASSO regression
Fit LASSO on the same feature set.
Purpose:
    sparse feature selection
    convex optimization model
    easier interpretation

11. Elastic net regression
Fit elastic net on the same feature set.
Purpose:
    combines LASSO sparsity with ridge stability
    usually better when predictors are correlated
    likely your best main model


# Validation / Analysis
17. Compare models
Compare:
    mean baseline
    previous glucose baseline
    ridge
    LASSO
    elastic net

18. Analyze selected predictors
Look at which variables elastic net/LASSO keeps:
    prior glucose
    insulin types
    meal effects
    exercise effects
    time-of-day effects

19. Interpret results
Explain what the selected predictors mean clinically and what the tradeoff is between:
    fewer predictors
    lower prediction error


# Concise version
Load and merge patient files
Transform raw date/time into Day and MinuteOfDay
Select pre-meal glucose rows as targets
Engineer lagged insulin, meal, exercise, and glucose features
Clean invalid or missing entries
Split data into past-for-training and future-for-testing
Fit baseline models
Fit ridge regression
Fit LASSO regression
Fit elastic net regression
Tune regularization parameters
Choose the best sparse model
Test on held-out data using MSE/RMSE
Interpret selected predictors and conclusions